\originalchapter{电势与连续电荷}\label{ch:potential}
\originalsection{功、势能与电势}
提升物体需要克服重力而增加势能, 移动电荷克服静电力也有相同的能量账本. 电场作功$W_{a\to b}=q\int_a^b\ee\cdot\dd\rr$等于势能减少量. 静电力是保守力, 因而定义每单位电荷的势能$V=U/q$后,
\[
V(b)-V(a)=-\int_a^b\ee\cdot\dd\rr,\qquad \ee=-\nabla V.
\]
势差单位为伏特(V), $1\,\mathrm V=1\,\mathrm{J/C}$. 电势的零点可任意选择, 同一配置中必须一致; 电场和可测势差不随整体加常数改变. 正电荷自发加速向低势, 负电荷向高势, 因为势能是$qV$, 不是单独$V$.

对有限的孤立电荷分布, 取无穷远势为零, 积分点电荷场可得$V=\kc q/r$. 多电荷势是标量相加, $V=\kc\sum_i q_i/|\rr-\rr_i|$. 组装点电荷系统时逐个加入电荷, 每一对只计一次, 所以$U=\sum_{i<j}\kc q_iq_j/r_{ij}$. 这不包括理想点电荷发散的自能.

等势面满足$V$恒定. 沿面移动时$\dd V=0$, 故$\ee$与面切向垂直. 等势面间距小表示势变化快、场较强; 两条场线间距和两等势面间距描述不同几何对象. $V$可以为零而场非零, 例如偶极子的中垂面.

静电力的功总满足$W_{\rm field}=-\Delta U$. 外力功却只有在动能不改变时才等于$\Delta U$; 一般功--能定理给$W_{\rm ext}=\Delta K+\Delta U$. “准静态搬运”就是缓慢且动能变化可忽略的理想过程, 不是任意快速推动电荷都得到同一外力功.

\originalsection{连续分布的积分}
对观测点$\rr$与源点$\rr'$分别标记, 则在源外
\[
V(\rr)=\kc\int\frac{\dd q}{|\rr-\rr'|},\qquad
\ee(\rr)=\kc\int\frac{\rr-\rr'}{|\rr-\rr'|^3}\dd q.
\]
分母是源到观测点的距离, 不是源点到坐标原点的距离. 势积分较简单, 但由势求场时仍须对观测坐标求导. 无穷长直线的势不能以无穷远为零, 应求两个有限半径之间的势差.

\begin{example}
均匀带电圆环半径$R$, 总电荷$Q$, 求轴上距环心$z$处的势与场.
\end{example}
\begin{solution}
所有小元到观测点的距离都是$\sqrt{R^2+z^2}$, 故$V=\kc Q/\sqrt{R^2+z^2}$. 环上相反小元的横向场抵消, 只剩轴向场$E_z=-\dd V/\dd z=\kc Qz/(R^2+z^2)^{3/2}$. $z=0$处场为零, 但势一般不为零; $|z|\gg R$时回到点电荷远场.
\end{solution}

均匀圆盘可拆成半径$s$、宽$\dd s$的圆环, $\dd q=2\pi\sigma s\dd s$. 在$z>0$的轴线上,
\[
E_z=2\pi\kc\sigma z\int_0^R\frac{s\dd s}{(s^2+z^2)^{3/2}}
=\frac\sigma{2\eps}\left(1-\frac z{\sqrt{R^2+z^2}}\right).
\]
当$R\to\infty$且保持$\sigma$固定时, 得到无限平面一侧的$E=\sigma/(2\eps)$; 当$z\gg R$时, 使用$Q=\pi R^2\sigma$回到点电荷场. 这两种极限固定的物理量不同, 不可混用.

\originalsection{偶极子的远场}
两电荷位于$\pm\vct d/2$, 且$r\gg d$, 将$1/|\rr\mp\vct d/2|$展开到一阶得到$1/r\pm(\vct d\cdot\uv r)/(2r^2)$. 两电荷单极项相消, 所以$V\simeq\kc\vct p\cdot\uv r/r^2$. 在沿$\vct p$的球坐标中, $V=\kc p\cos\theta/r^2$, 从径向和切向变化率求得
\[
E_r=\frac{2\kc p\cos\theta}{r^3},\qquad E_\theta=\frac{\kc p\sin\theta}{r^3}.
\]
这里保留的是远场首项, 不能把它当成两电荷附近的精确公式. 偶极场按$r^{-3}$衰减, 比有净电荷的$r^{-2}$场更快.


展开可从$|\rr\mp\vct d/2|^{-1}=r^{-1}[1\mp(\uv r\cdot\vct d)/r+d^2/(4r^2)]^{-1/2}$和$(1+u)^{-1/2}=1-u/2+O(u^2)$得到. 沿极角移动$\dd\theta$的弧长为$r\dd\theta$, 所以切向场$E_\theta=-(1/r)\partial_\theta V$, 代入偶极势即得到上式.

固定的微小偶极在光滑静电外场中, 两端力差为$\vct F\simeq q(\vct d\cdot\nabla)\ee=(\vct p\cdot\nabla)\ee$. 静电旋度为零, 所以也等于$\nabla(\vct p\cdot\ee)$. 例如在对称轴上$\vct p=p\uv z$, $E_z$随$z$变化, 轴向力$F_z=p\dd E_z/\dd z$. 对局部轴上表达$E_z(z)=E_*+\beta z$, 对齐偶极受$p\beta$, 翻转偶极受反向力. 这里不把只有$E_z(z)$的全空间场冒充无电荷场.
